% GENERATED FILE. Edit canonical lessons, then run npm run generate:books.
% canonical-source-hash: fnv1a64:e56200f51d2c6164
% canonical-lessons: FA-C102-zanjabil, FA-C102-felfel, FA-C102-ghahve, FA-C102-sharbat, FA-C102-dasht

\chapter{Ginger, Pepper, Coffee, Sweet drink, Plain}
\label{ch:fa-ginger-pepper-coffee-sweet-drink-plain}
\hlchaptermodality{102}

\begin{chapteropening}
\textbf{By the end of this chapter:} \emph{I can say ginger, pepper, coffee, a sweet drink and a plain.}

Complete the last lesson of chapter 102: I can say ginger, pepper, coffee, a sweet drink and a plain.
\end{chapteropening}

\section[\texorpdfstring{\fa{زنجبیل} (zanjabil)}{زنجبیل (zanjabil)}]{\fa{زنجبیل} (zanjabil) — ginger}

\label{lesson:FA-C102-zanjabil}



\begin{quote}
Before the new one: say the Persian for an aubergine, then the Persian for garlic.
\end{quote}

\subsection*{You'll want to know: \fa{زنجبیل}}
\textbf{\fa{زنجبیل}} — \emph{zanjabil} — \textquotedblleft{}ginger\textquotedblright{}.

\begin{cousinweb}
\emph{zanjabil} is borrowed from Arabic \emph{zanjabīl}, \textquotedblleft{}ginger\textquotedblright{}, which came through Middle Persian \emph{singibēr} from an Indian word (EtymDB). English \emph{ginger} goes back to the same Indian word through Greek and Latin.
\end{cousinweb}

\begin{grammarlens}[title={What you've built}]
One more word for this chapter.
\end{grammarlens}

\subsection*{Your turn}
\begin{itemize}
\raggedright
  \item \emph{Say it:} \emph{zanjabil}
  \item \emph{Say it:} \emph{zanjabil}, once more
  \item \emph{Say it:} say \emph{zanjabil}
  \item \emph{Recall:} say the Persian for an aubergine, then the Persian for garlic, then say \emph{zanjabil} again
\end{itemize}

\begin{tcolorbox}[breakable,colback=teal!4,colframe=teal!35!black,title={Before you move on}]
What does \fa{زنجبیل} mean? (\textquotedblleft{}Ginger\textquotedblright{}.) Say it once more.
\end{tcolorbox}
\section[\texorpdfstring{\fa{فلفل} (felfel)}{فلفل (felfel)}]{\fa{فلفل} (felfel) — pepper}

\label{lesson:FA-C102-felfel}



\begin{quote}
Before the new one: say the Persian for garlic, then the Persian for ginger.
\end{quote}

\subsection*{You'll want to know: \fa{فلفل}}
\textbf{\fa{فلفل}} — \emph{felfel} — \textquotedblleft{}pepper\textquotedblright{}.

\begin{cousinweb}
\emph{felfel} is borrowed from Arabic \emph{filfil}, \textquotedblleft{}pepper\textquotedblright{}, which is usually traced through Persian \emph{pilpil} to Sanskrit \emph{pippalī}, \textquotedblleft{}long pepper\textquotedblright{}. English \emph{pepper} comes from the same Indian word through Greek and Latin.
\end{cousinweb}

\begin{grammarlens}[title={What you've built}]
One more word for this chapter.
\end{grammarlens}

\subsection*{Your turn}
\begin{itemize}
\raggedright
  \item \emph{Say it:} \emph{felfel}
  \item \emph{Say it:} \emph{felfel}, once more
  \item \emph{Say it:} say \emph{felfel}
  \item \emph{Recall:} say the Persian for garlic, then the Persian for ginger, then say \emph{felfel} again
\end{itemize}

\begin{tcolorbox}[breakable,colback=teal!4,colframe=teal!35!black,title={Before you move on}]
What does \fa{فلفل} mean? (\textquotedblleft{}Pepper\textquotedblright{}.) Say it once more.
\end{tcolorbox}
\section[\texorpdfstring{\fa{قهوه} (ghahve)}{قهوه (ghahve)}]{\fa{قهوه} (ghahve) — coffee}

\label{lesson:FA-C102-ghahve}



\begin{quote}
Before the new one: say the Persian for ginger, then the Persian for pepper.
\end{quote}

\subsection*{You'll want to know: \fa{قهوه}}
\textbf{\fa{قهوه}} — \emph{ghahve} — \textquotedblleft{}coffee\textquotedblright{}.

\begin{cousinweb}
\emph{ghahve} is borrowed from Arabic \emph{qahwa}, \textquotedblleft{}coffee, a brew\textquotedblright{} (EtymDB). Its Arabic root is \emph{q-h-w}.
\end{cousinweb}

\begin{grammarlens}[title={What you've built}]
One more word for this chapter.
\end{grammarlens}

\subsection*{Your turn}
\begin{itemize}
\raggedright
  \item \emph{Say it:} \emph{ghahve}
  \item \emph{Say it:} \emph{ghahve}, once more
  \item \emph{Say it:} say \emph{ghahve}
  \item \emph{Recall:} say the Persian for ginger, then the Persian for pepper, then say \emph{ghahve} again
\end{itemize}

\begin{tcolorbox}[breakable,colback=teal!4,colframe=teal!35!black,title={Before you move on}]
What does \fa{قهوه} mean? (\textquotedblleft{}Coffee\textquotedblright{}.) Say it once more.
\end{tcolorbox}
\section[\texorpdfstring{\fa{شربت} (sharbat)}{شربت (sharbat)}]{\fa{شربت} (sharbat) — a sweet drink}

\label{lesson:FA-C102-sharbat}



\begin{quote}
Before the new one: say the Persian for pepper, then the Persian for coffee.
\end{quote}

\subsection*{You'll want to know: \fa{شربت}}
\textbf{\fa{شربت}} — \emph{sharbat} — \textquotedblleft{}a sweet drink\textquotedblright{}.

\begin{cousinweb}
\emph{sharbat} is borrowed from Arabic \emph{šarba}, \textquotedblleft{}a drink, a draught\textquotedblright{}. Its Arabic root is \emph{š-r-b} (Lane). English \emph{sherbet} comes from the same word through Turkish.
\end{cousinweb}

\begin{grammarlens}[title={What you've built}]
One more word for this chapter.
\end{grammarlens}

\subsection*{Your turn}
\begin{itemize}
\raggedright
  \item \emph{Say it:} \emph{sharbat}
  \item \emph{Say it:} \emph{sharbat}, once more
  \item \emph{Say it:} say \emph{sharbat}
  \item \emph{Recall:} say the Persian for pepper, then the Persian for coffee, then say \emph{sharbat} again
\end{itemize}

\begin{tcolorbox}[breakable,colback=teal!4,colframe=teal!35!black,title={Before you move on}]
What does \fa{شربت} mean? (\textquotedblleft{}A sweet drink\textquotedblright{}.) Say it once more.
\end{tcolorbox}
\section[\texorpdfstring{\fa{دشت} (dasht)}{دشت (dasht)}]{\fa{دشت} (dasht) — a plain}

\label{lesson:FA-C102-dasht}



\begin{quote}
Before the new one: say the Persian for coffee, then the Persian for a sweet drink.
\end{quote}

\subsection*{You'll want to know: \fa{دشت}}
\textbf{\fa{دشت}} — \emph{dasht} — \textquotedblleft{}a plain\textquotedblright{}.

\begin{cousinweb}
\emph{dasht} is a native Persian word, inherited rather than borrowed; it is usually traced to Middle Persian \emph{dašt}, \textquotedblleft{}plain\textquotedblright{}. (EtymDB links it to the Indo-European word for \textquotedblleft{}right hand\textquotedblright{}, but that link is not settled.)
\end{cousinweb}

\begin{grammarlens}[title={What you've built}]
One more word for this chapter.
\end{grammarlens}

\subsection*{Your turn}
\begin{itemize}
\raggedright
  \item \emph{Say it:} \emph{dasht}
  \item \emph{Say it:} \emph{dasht}, once more
  \item \emph{Say it:} say \emph{dasht}
  \item \emph{Recall:} say the Persian for coffee, then the Persian for a sweet drink, then say \emph{dasht} again
\end{itemize}

\begin{tcolorbox}[breakable,colback=teal!4,colframe=teal!35!black,title={Before you move on}]
What does \fa{دشت} mean? (\textquotedblleft{}A plain\textquotedblright{}.) Say it once more.
\end{tcolorbox}
